\scanpage{36}
\mspage{3}
qu'on peut interpréter aussi de la façon suivante : La famille \inserted{$(\Omega_i)_{i\in I}$} des trois ensembles $\Omega_i$ à deux éléments définit un cube -- et on a une orientation de ce cube.

\uline{NB} Les $\Omega_i$ et $\Phi_i$ se déterminent mutuellement par
\[
\left\{\begin{aligned}
 \Omega_i&\simeq\Phi_i\wedge\bigwedge_{\alpha\in I\setminus\{i\}}(I\setminus\{\alpha\})\\
 \Phi_i&\simeq\Omega_i\wedge\bigwedge_{\alpha\in I\setminus\{i\}}(I-\{\alpha\})
\end{aligned}\right.\tag*{(13)}
\]
% slanted
\margin{\unsure{itou pour|idem pour} $\Omega_i$, $\Phi_i$ \deleted{$=$} \inserted{et} $A_i$, puisque
$\begin{aligned}
 A_i&\simeq\deleted{\unsure{\Omega_i\wedge|\bigwedge}}\,(I\setminus\{i\})\deleted{\unsure{S_i}}\\
 &\simeq\Phi_i\wedge\bigwedge_{\alpha\in I}(I-\alpha)\\
 &\simeq\Phi_i\wedge\operatorname{or}(I)
 \end{aligned}$\\ cf (9))}
qui donnent des $\mathrm{isom}^{\mathrm s}$ canoniques
\[
\left\{\begin{aligned}
 \bigwedge_{i\in I}\Omega_i&\simeq\bigwedge_{i\in I}\Phi_i\\
 \text{et}\qquad
 \underbrace{\left(\bigwedge_{i\in I}\Omega_i\right)\wedge\operatorname{or}\,I}
 _{\substack{\text{orientations \unsure{de|avec} l'}\\\text{\unsure{oct.|oct} comb. défini}\\\text{par }(\Omega_i)_{i\in I}}}
 &\simeq
 \underbrace{\left(\bigwedge_{i\in I}\Phi_i\right)\wedge\operatorname{or}\,I}
 _{\substack{\text{orientations de l'oct.}\\\text{comb. défini par}\\(\Phi_i)_{i\in I}.}}
\end{aligned}\right.\tag*{(14)}
\]

En résumé

\uline{Proposition} Considérons la carte combinatoire $\unsure{\mathcal C|\mathscr C}$ (avec trois sommets, six arêtes, quatre faces qui sont des triangles) correspondant à un ``arrangement (combinatoire) de \deleted{quatre} \inserted{trois} ps-droites \inserted{$D_i$ $(i\in I)$} du plan projectif''. Soit $K$ son
% slanted
\margin{octaèdre gauche}
% archive sheet number: 35
